\chapter{Drug Overdose}

\section*{Definition and Overview}
A drug overdose occurs when a person takes more of a drug than the body can handle, whether it is a prescription medicine, an over-the-counter product, alcohol, or an illicit drug. Overdoses can be accidental, deliberate, or the result of drug misuse, and they range from mild effects to life-threatening emergencies. The signs depend on the drug involved: opioids and sedatives slow breathing, stimulants overheat and strain the heart, and paracetamol causes delayed liver damage. An overdose is a medical emergency, and immediate help can save a life. This chapter explains the signs of overdose and what to do.

\section*{Epidemiology}
Drug overdoses cause a significant number of deaths and hospitalisations in Australia each year. Opioids, including prescription opioids and heroin, are the leading cause of overdose deaths, and accidental overdose is a growing concern. Paracetamol overdose is common and potentially fatal, and alcohol overdose occurs in binge drinkers. Overdose risk is higher with mixing drugs, taking drugs after a period of abstinence, and using drugs alone. Naloxone, a medicine that reverses opioid overdose, is now available to people who use opioids and their families.

\section*{Aetiology and Pathophysiology}
\begin{itemize}
    \item \textbf{Opioids:} slow the brain and breathing, causing respiratory depression, coma, and death
    \item \textbf{Sedatives and alcohol:} depress the central nervous system; combining them is especially dangerous
    \item \textbf{Stimulants:} raise heart rate and blood pressure, causing heart attack, stroke, and hyperthermia
    \item \textbf{Paracetamol:} causes liver damage that develops over days
    \item \textbf{Other medicines:} antidepressants, antipsychotics, and other drugs have specific toxicity
    \item \textbf{Mechanism:} the drug overwhelms the body's systems, disrupting breathing, heart function, or organs
\end{itemize}

\section*{Clinical Features}
\begin{itemize}
    \item Opioid overdose: pinpoint pupils, slow or absent breathing, deep sleep, and unresponsiveness
    \item Alcohol or sedative overdose: confusion, vomiting, slow breathing, and unconsciousness
    \item Stimulant overdose: agitation, fast heart rate, high blood pressure, chest pain, and seizures
    \item Paracetamol overdose: few symptoms initially, then nausea, vomiting, and liver failure over days
    \item Unconsciousness, seizures, and collapse in severe cases
    \item Blue lips or skin from lack of oxygen
\end{itemize}

\section*{Differential Diagnosis}
\begin{itemize}
    \item Overdose of different drug classes, which need different emergency treatment
    \item Overdose versus other causes of collapse, such as stroke, seizures, and low blood sugar
    \item Accidental versus deliberate overdose, which affects aftercare
    \item Poisoning from other substances, including household products
\end{itemize}

\section*{When to Seek Medical Care}
Any suspected overdose requires immediate emergency care. Do not wait for symptoms to improve. If someone is unresponsive, not breathing normally, or seriously unwell after taking drugs, call for an ambulance.

\textbf{Call 000 (or local emergency services) for:}
\begin{itemize}
    \item Unresponsiveness, difficulty waking, or no breathing
    \item Slow, shallow, or irregular breathing
    \item Seizures, agitation, or collapse
    \item Chest pain, fast heart rate, or difficulty breathing after stimulants
    \item Any suspected overdose, even if the person seems to improve
\end{itemize}

\section*{Diagnosis and Investigations}
\begin{itemize}
    \item \textbf{Assessment of airway, breathing, and circulation:} the first priority
    \item \textbf{History:} what was taken, how much, and when
    \item \textbf{Blood tests:} paracetamol and other drug levels, electrolytes, and organ function
    \item \textbf{Toxicology screens:} identify drugs present
    \item \textbf{ECG and monitoring:} for heart effects
    \item \textbf{Imaging:} in some cases, such as suspected ingestion of drug packets
\end{itemize}

\section*{Management}
\begin{itemize}
    \item \textbf{Emergency care:} airway support, oxygen, and breathing support in hospital
    \item \textbf{Naloxone:} reverses opioid overdose and can be given by bystanders or paramedics
    \item \textbf{Activated charcoal:} may be given early for some ingestions
    \item \textbf{Antidotes:} specific treatments for particular drugs
    \item \textbf{Paracetamol antidote:} N-acetylcysteine prevents liver damage
    \item \textbf{Supportive care:} fluids, temperature control, and monitoring
    \item \textbf{Mental health assessment:} after deliberate overdose
    \item \textbf{Follow-up:} addiction treatment and support to prevent recurrence
\end{itemize}

\section*{Complications}
\begin{itemize}
    \item Death from respiratory failure or heart problems
    \item Brain injury from oxygen deprivation
    \item Liver failure from paracetamol
    \item Kidney damage and muscle breakdown
    \item Aspiration of vomit into the lungs
    \item Psychological and social consequences
\end{itemize}

\section*{Prognosis and Outcomes}
The outcome depends on the drug, the amount, the speed of treatment, and the person's health. With prompt emergency care, many overdoses are fully reversible, and naloxone has saved many lives from opioid overdose. Paracetamol overdose is highly treatable when treated early. Survivors of overdose benefit from follow-up, including mental health support and addiction treatment, which reduce the risk of recurrence.

\section*{Prevention and Self-Care}
\begin{itemize}
    \item Take medicines exactly as prescribed and store them safely
    \item Never mix drugs, including alcohol, without medical advice
    \item Avoid using drugs alone
    \item Carry naloxone if you or someone you know uses opioids
    \item Keep medicines out of reach of children
    \item Follow dosing instructions for over-the-counter products
    \item Seek help for addiction and mental health problems
    \item Call an ambulance immediately if an overdose is suspected
\end{itemize}

\section*{Sources and Further Reading}
The following reputable sources provide further information for healthcare students and patients seeking reliable, current advice about drug overdose.

\begin{itemize}
    \item Alcohol and Drug Foundation. \textit{Overdose: prevention and response.} https://adf.org.au/
    \item Healthdirect Australia. \textit{Drug overdose: emergency information.} https://www.healthdirect.gov.au/
    \item Penington Institute. \textit{Overdose prevention and naloxone.} https://www.penington.org.au/
    \item National Coronial Information System. \textit{Drug-induced deaths in Australia.} https://www.ncis.org.au/
    \item MedlinePlus. \textit{Drug overdose.} https://medlineplus.gov/
    \item Australian Resuscitation Council. \textit{First aid for overdose.} https://resus.org.au/
\end{itemize}
